\title{FlashAttention-3: Fast and Accurate Attention with Asynchrony and Low-precision}

\begin{abstract}
As the foundational component of the Transformer architecture, attention serves as the principal constraint in scaling large language models and in managing extended context sequences. \textsc{FlashAttention} introduced a GPU-optimized method by reducing memory access overhead; nonetheless, this approach does not fully leverage the advancements available in state-of-the-art hardware. Notably, \textsc{FlashAttention-2} achieves only 35\% throughput utilization on the H100 GPU. To address this gap, we present three novel strategies for accelerating attention on Hopper GPUs: harnessing the asynchrony between Tensor Cores and TMA to (1) interleave computation and memory transfers through warp-level specialization, (2) integrate block-wise matrix multiplication with softmax computations, and (3) apply block-level quantization with asynchronous processing to utilize native FP8 low-precision capabilities. Our proposed solution, \textsc{FlashAttention-3}, attains a 1.5–2.0$\times$ speedup on H100 GPUs, achieving performance of up to 740 TFLOPs/s with FP16 (corresponding to 75\% hardware utilization) and nearly 1.2 PFLOPs/s with FP8. We further show that FP8 \textsc{FlashAttention-3} yields numerical errors 2.6$\times$ lower than those observed in a standard FP8 attention implementation.
\end{abstract}

\section{Introduction}
\label{sec:intro}

Within the Transformer architecture~\citep{vaswani2017attention}, the computation of self-attention forms the dominant computational challenge, as calculating the attention scores between queries and keys exhibits quadratic complexity with respect to sequence length. Enabling attention to handle substantially longer contexts could facilitate a range of novel capabilities, such as reasoning across several extended documents~\citep{guo2021longt5,shaham2022scrolls,peng2023yarn} and analyzing files within expansive code repositories~\citep{roziere2023code, li2023starcoder}; accommodating additional data modalities, for example, detailed image data~\citep{chen2022scaling}, audio streams~\citep{gulati2020conformer}, and video sequences~\citep{ho2022video}; or unlocking new use cases, such as processing extensive user interaction histories~\citep{sun2019bert4rec} and managing agent tasks with distant objectives~\citep{yao2022react}. These prospects have motivated substantial efforts to accelerate attention mechanisms in the setting of long contexts, either through approximate methods~\citep{katharopoulos2020transformers,choromanski2020rethinking,
tay2020efficient}, advances in software efficiency (\citep{rabe2021self,dao2022flashattention,kwon2023efficient}), or the adoption of alternative architectural paradigms~\citep{peng2023rwkv,sun2023retentive,gu2023mamba}.

This paper extends the research by \citet{dao2022flashattention}, who pioneered exact-attention algorithms that consciously incorporate GPU execution paradigms and hardware capabilities into the algorithmic design process. Dao et al. in \citep{dao2022flashattention} presented \textsc{FlashAttention}, an innovative tiling method which parallelizes attention computation by merging all attention steps into one GPU kernel, thereby avoiding unnecessary transfers to and from global memory.

Later developments by \citet{dao2023flashattention2} led to \textsc{FlashAttention-2}, which further reorganized the computation by distributing work not only across the batch and head dimensions, but also the sequence length. This version carries out the forward computation over blocks of the key and value matrices, with the intention of enhancing GPU utilization and balancing the work load. Despite these changes, we note that \textsc{FlashAttention-2} does not attain optimal throughput on recent hardware; for instance, its utilization on the Hopper H100 GPU can be as low as 35\%, compared to the 80–90\% achieved by highly-tuned GEMM kernels. This efficiency gap may be explained in part by lower-level details, such as reliance on Ampere-specific instructions instead of those specialized for Hopper when employing Tensor Cores. Other projects like ThunkerKitten~\citep{spector2024thunder} and cuDNN 9~\citep{cudnn9} have demonstrated that leveraging Hopper-optimized instructions and adopting tile-oriented approaches can both accelerate attention calculations and streamline code complexity.

At a deeper level, the \textsc{FlashAttention-2} algorithm is built on a basic synchronous computation framework and does not inherently exploit asynchronous execution or reduced-precision arithmetic. Hardware advances, however, have introduced specialized units to optimize performance for critical machine learning operations: for instance, dedicated components like Tensor Cores accelerate matrix multiplication, while others such as the Tensor Memory Accelerator (TMA) handle memory transfers, operating independently from general-purpose CUDA cores that execute the remaining computational tasks (logic, integer, or standard floating point arithmetic). Furthermore, support for low-precision data types—ranging from FP16 in Pascal (2017), through BF16 in Ampere (2020), to the most recent FP8 in Hopper and FP4 in Blackwell—substantially increases computational throughput without raising power consumption or silicon area, following a proven progression in hardware capabilities. A discussion of such enhancements specific to Hopper appears in \cref{subsec:hardware}.  
The engineering problem thus centers on how to rearchitect \textsc{FlashAttention-2} so it can harness these hardware capabilities: exploiting asynchrony demands simultaneous execution of operations like matrix multiplication and softmax—despite inherent data dependencies between them—while employing low-precision computation mandates sophisticated handling to prevent excessive quantization errors, which is particularly pertinent for capturing outlier activations in large language models~\citep{dettmers2208llm, sun2024massive}.

With this goal in mind, we introduce \textsc{FlashAttention-3}, which integrates three novel concepts to enhance efficiency on cutting-edge GPU architectures:\footnote{Our discussion is framed around NVIDIA's Hopper architecture. Nevertheless, the presented algorithm applies to any GPU architecture that offers adequate support for asynchronous operations and low-precision computation.}

\begin{enumerate}

\item \textbf{Producer-Consumer asynchrony:} We introduce a warp-specialized software pipelining approach that leverages the asynchronous execution characteristics of data movement operations and Tensor Cores by assigning producer and consumer roles to distinct warps. This separation extends the algorithm’s potential to effectively obscure both memory access and instruction issue latencies.

\item \textbf{Hiding softmax under asynchronous block-wise GEMMs:} We concurrently execute the relatively lower-throughput softmax-related operations, such as floating point multiply-adds and exponentiations, alongside asynchronous WGMMA operations for GEMM. In order to facilitate this concurrency, we revise the \textsc{FlashAttention-2} methodology to avoid certain inherent sequential dependencies between the softmax and GEMM computations. For instance, in our 2-stage algorithmic variant, softmax is processed on a given score matrix block at the same time as WGMMA asynchronously calculates the succeeding block.

\item \textbf{Hardware-accelerated low-precision GEMM:} We update the forward pass process to utilize FP8 Tensor Cores for GEMM tasks, which results in a nearly two-fold increase in achieved TFLOPs/s. Addressing this necessitates accommodating the distinct memory layout constraints imposed by WGMMA regarding the arrangement of FP32 accumulator and FP8 operand matrix blocks. Techniques such as block quantization and incoherent processing are implemented to reduce the degradation in computational accuracy typically associated with the transition to FP8 precision.
\end{enumerate}

To empirically assess the performance of our approach, we conduct benchmarks of \textsc{FlashAttention-3} on the H100 SXM5 GPU across various settings. The results demonstrate that (1) FP16 provides a speedup of 1.5-2.0$\times$ over \textsc{FlashAttention-2} during the forward pass—achieving up to 740 TFLOPs/s—and a 1.5-1.75$\times$ improvement in the backward pass; (2) FP8 attains nearly 1.2 PFLOPs/s; and (3) for larger sequence lengths, FP16 not only surpasses alternatives, but FP8 performs on par with or exceeds\footnote{Specifically, with a head dimension of 64, \textsc{FlashAttention-3} FP8 has an advantage, while for head dimensions of 128 and 256, it matches performance in non-causal masking scenarios but lags behind in those with causal masking.} a leading attention implementation from NVIDIA's cuDNN library. In addition, our evaluation confirms that FP16 \textsc{FlashAttention-3} preserves the numerical error characteristics of \textsc{FlashAttention-2}, outperforming conventional attention computations, partly because intermediate quantities (such as softmax rescaling) remain in FP32. Notably, FP8 \textsc{FlashAttention-3}, when combined with block quantization and incoherent processing, exhibits 2.6$\times$ greater accuracy compared to conventional attention employing per-tensor quantization, especially in cases involving outlier features.

We have released \textsc{FlashAttention-3} as open source under a permissive license\footnote{\textsc{FlashAttention-3}
  is available at \texttt{https://github.com/Dao-AILab/flash-attention}}, and intend to incorporate it into both PyTorch and Hugging Face ecosystems to maximize its accessibility and impact among researchers and developers.

\section{Background: Multi-Head Attention and GPU Characteristics}
\label{sec:background}

\subsection{Multi-Head Attention}
\label{subsec:multi_head_attn}

Let $\mathbf{Q}, \mathbf{K}, \mathbf{V} \in \mathbb{R}^{N \times d}$ be the query, key and value input sequences associated to a single head, where $N$ is the sequence length and $d$ is the head dimension. Then the attention output $\mathbf{O}$ is computed as:

\begin{equation*}
  \mathbf{S} = \alpha \mathbf{Q} \mathbf{K}^\top \in \mathbb{R}^{N \times N}, \quad \mathbf{P} = \mathrm{softmax}(\mathbf{S}) \in \mathbb{R}^{N \times N}, \quad \mathbf{O} = \mathbf{P}\mathbf{V} \in \mathbb{R}^{N \times d},
\end{equation*}

Here, $\mathrm{softmax}$ operates along the rows, and the scaling coefficient $\alpha$ is generally chosen to be $1/\sqrt{d}$. To mitigate issues of numerical instability arising from exponentiation, it is standard practice to subtract $\mathrm{rowmax}(\mathbf{S})$ from $\mathbf{S}$. In the context of multi-head attention (MHA), each attention head utilizes independent query, key, and value projection matrices. This process is performed in parallel over various heads and batches, resulting in the complete output tensor.

Now let $\phi$ be a scalar loss function and let $\mathbf{d}(-) = \partial \phi / \partial (-)$ be notation for the gradient. Given the output gradient $\mathbf{dO} \in \mathbb{R}^{N \times d}$, we compute $\mathbf{dQ}$, $\mathbf{dK}$, and $\mathbf{dV}$ according to the chain rule as follows:

\begin{align*}
  \mathbf{dV} &= \mathbf{P}^\top \mathbf{dO} \in \mathbb{R}^{N \times d} \\
  \mathbf{dP} &= \mathbf{dO} \mathbf{V}^\top \in \mathbb{R}^{N \times N} \\
  \mathbf{dS} &= \mathrm{dsoftmax} (\mathbf{dP}) \in \mathbb{R}^{N \times N} \\
  \mathbf{dQ} &= \alpha \mathbf{dS} \mathbf{K} \in \mathbb{R}^{N \times d} \\
  \mathbf{dK} &= \alpha \mathbf{dS}^\top \mathbf{Q} \in \mathbb{R}^{N \times d},
\end{align*}

In this context, $\mathbf{d}s = (\mathrm{diag}(p) - p p^\top)\mathbf{d}p$ holds when $p = \mathrm{softmax}(s)$, where $p$ is represented as a function of the vector $s$. The notation $\mathrm{dsoftmax}(\mathbf{dP})$ refers to the application of this formula to each row individually. Importantly, this step also lends itself well to parallelization over both the number of heads and batches during the backward pass of MHA.

\subsection{GPU hardware characteristics and execution model}
\label{subsec:hardware}

We outline the components of the GPU execution model pertinent to \textsc{FlashAttention-3}, concentrating specifically on the NVIDIA Hopper architecture as a representative example of this framework.

\paragraph{Memory hierarchy:} The GPU features a hierarchical structure of memory locations, with bandwidth generally decreasing as capacity increases (\cref{tab:gpu-hierarchy})\footnote{\citet{luo2024benchmarking} indicates that shared memory achieves a bandwidth of 128 bytes per clock cycle per SM; this figure is scaled by the number of SMs, 132, and the boost clock frequency of 1830 MHz.}. At the highest level, global memory (GMEM)—also known as HBM—serves as off-chip DRAM and is shared among all streaming multiprocessors (SMs). The L2 cache, situated on-chip, automatically caches data from GMEM. Each SM further possesses its own small, on-chip, banked, programmer-controlled shared memory (SMEM). The memory closest to computation is the register file located within each SM.

\paragraph{Thread hierarchy:} In the GPU programming paradigm, execution units are systematically arranged into a thread hierarchy. This structure starts at individual threads, which are organized into warps consisting of 32 threads each. Four sequential warps are combined to form a warpgoup. Warpgroups are then aggregated into threadblocks (also known as cooperative thread arrays or CTAs). In the Hopper architecture, multiple threadblocks can be further grouped into threadblock clusters, and finally, all of these are contained within grids at the highest level.

There exists a strong interconnection between these two hierarchies. All threads belonging to a single CTA are assigned to execute concurrently on the same SM, while multiple CTAs within a cluster are concurrently dispatched to a single GPC. Every thread within a CTA has shared access to SMEM, in contrast to RMEM, where each thread has exclusive use of up to 256 private registers.

\paragraph{Asynchrony and warp-specialization:}

GPUs achieve high throughput by leveraging parallelism and asynchronous execution to mask delays arising from memory access and computation. Hopper introduces the Tensor Memory Accelerator (TMA), a specialized hardware component dedicated to enabling asynchronous memory transfers between GMEM and SMEM \cite[\S7.29]{cuda}. Additionally, departing from earlier generations like Ampere, the Tensor Core in Hopper—accessible through the warpgroup-wide WGMMA instruction \cite[\S9.7.14]{ptx}—also operates asynchronously and is capable of fetching its operands straight from shared memory.

Enabling asynchrony through hardware facilitates the design of warp-specialized kernels, where the warps within a CTA are assigned exclusively to either producer or consumer roles, engaging solely in data transfer or computation, respectively. This strategy generally enhances the compiler’s capability to arrange instructions more effectively \citep{warp-specialization-2011}. Moreover, Hopper architecture introduces support for dynamically redistributing registers among warpgroups using the \verb|setmaxnreg| operation \cite[\S9.7.17.1]{ptx}, allowing warps performing MMAs to access a greater portion of RMEM in comparison to those limited to executing TMA instructions, which necessitate only a single active thread.

\paragraph{Low-precision number formats:}
\label{sec:low-precision-gpu}
Recent GPU architectures are equipped with dedicated hardware components designed to enhance the speed of low-precision arithmetic. As an illustration, utilizing the WGMMA instruction enables FP8 Tensor Cores on Hopper to achieve double the throughput per SM relative to FP16 or BF16 performance.

Nonetheless, to utilize FP8 WGMMA accurately, it is essential to recognize the specific memory layout requirements imposed on its operands. In the context of a GEMM operation computing $A \times B^{\top}$, with $A$ as an $M\times K$ matrix and $B$ as an $N\times K$ matrix, we classify an operand as \emph{mn-major} when data is stored contiguously along the outermost $M$ or $N$ dimension, and as \emph{k-major} if the data is contiguous along the inner $K$ dimension. While FP16 WGMMA allows both mn-major and k-major operand layouts when accessing data in SMEM, FP8 WGMMA restricts support solely to the k-major arrangement. Additionally, in contexts such as attention mechanisms, where there is a need to chain multiple GEMM computations within a single kernel launch, incompatibilities between the memory layouts of FP32 accumulators and FP8 operands complicate the chaining of dependent FP8 WGMMA operations.

Within the scope of attention mechanisms, these layout constraints necessitate specific adjustments to how an FP8 algorithm is constructed, as outlined in \cref{sec:algofp8}.

\subsection{Standard Attention and Flash Attention} In accordance with \citet{dao2022flashattention}, we refer to \textbf{standard attention} as the GPU-based method for computing attention that stores the intermediate results, specifically the matrices $\mathbf{S}$ and $\mathbf{P}$, in HBM. The core innovation of \textsc{FlashAttention} involves the use of a local softmax computation, which eliminates the costly intermediate storage by integrating the attention computation within a single kernel call. This localized softmax operation is implemented in lines \ref{code-ws:softmax_start}-\ref{code-ws:softmax_end} of the consumer mainloop described in \cref{alg:flash3_wgmma_ws_only}, together with the block-wise rescaling of the output matrix $\mathbf{O}$. A straightforward proof demonstrating that this approach correctly computes $\mathbf{O}$ is available in \cite[\S 2.3.1]{dao2023flashattention2}.

\section{FlashAttention-3: Algorithm}
\label{sec:algo}

In this section, we outline the \textsc{FlashAttention-3} algorithm, concentrating on the forward pass; details of the backward pass can be found in~\cref{sec:algo_ws_bwd}. Initially, we illustrate the integration of warp-specialization alongside a circular SMEM buffer within the foundational \textsc{FlashAttention-2} algorithm. Next, we detail the approach to leverage WGMMA’s asynchronous capabilities to construct a two-stage pipeline that overlaps GEMM and softmax computations. Lastly, we discuss the required adjustments for FP8 support, addressing both compliance with data layout and maintenance of precision through block quantization and incoherent execution strategies.

\subsection{Producer-Consumer asynchrony through warp-specialization and pingpong scheduling}
\label{sec:algo_ws}

\paragraph{Warp-specialization}
Similar to \textsc{FlashAttention-2}, the forward computation in \textsc{FlashAttention-3} can be efficiently parallelized across the batch dimension, the number of attention heads, and the length of the query sequence. Consequently, it is adequate to analyze the algorithm at the CTA granularity, where each CTA processes a specific submatrix $\mathbf{Q}_i$ from the query matrix and produces the corresponding output tile $\mathbf{O}_i$. For clarity, we initially present the warp-specialization method using a circular SMEM buffer, without incorporating GEMM-softmax overlapping at this stage. Denote the head dimension by $d$ and the sequence length by $N$, and choose a query block size $B_r$ so that $\mathbf{Q}$ is segmented into $T_r = \lceil \frac{N}{B_r} \rceil$ blocks labeled $\mathbf{Q}_1, .., \mathbf{Q}_{T_r}$.

\begin{algorithm}[H]
    \caption{\small\label{alg:flash3_wgmma_ws_only}\textsc{FlashAttention-3} forward pass \textbf{absent} intra-consumer overlap -- CTA view}
    \begin{algorithmic}[1]
\REQUIRE Input matrices $\mathbf{Q}_i \in \mathbb{R}^{B_r \times d}$, $\mathbf{K}, \mathbf{V} \in \mathbb{R}^{N \times d}$ situated in HBM, specified key block size $B_c$ such that $T_c = \lceil \frac{N}{B_c} \rceil$.
\STATE Set up a pipeline controller for synchronizing via barriers with an $s$-stage circular shared memory buffer.
\IF {assigned to the producer warpgroup}
\STATE Release a fixed quantity of registers.
\STATE Schedule a transfer of $\mathbf{Q}_i$ from HBM to shared memory.
\STATE After the transfer, indicate completion to consumer threads for $\mathbf{Q}_i$'s availability.
\FOR{$0 \le j < T_c$}
    \STATE Await consumption of the $(j\,\%\,s)$th pipeline buffer stage.
    \STATE Initiate loads of $\mathbf{K}_j, \mathbf{V}_j$ from HBM into the shared memory buffer at stage $(j\,\%\,s)$.
    \STATE Following the loads, inform consumers of $\mathbf{K}_j, \mathbf{V}_j$ readiness.
\ENDFOR
\ELSE \STATE Allocate a predetermined number of registers, dependent on the count of consumer warps.
\STATE Locally, set initial values: $\mathbf{O}_i = (0) \in \mathbb{R}^{B_r \times d}$; $\ell_i = 0$, $m_i = -\infty$ in $\mathbb{R}^{B_r}$.
\STATE Block until $\mathbf{Q}_i$ arrives in shared memory.
\FOR{$0 \le j < T_c$}
\STATE Wait until $\mathbf{K}_j$ appears in shared memory.
\STATE Execute matrix multiplication: $\mathbf{S}_i^{(j)} = \mathbf{Q}_i \mathbf{K}_j^T$ (SS-GEMM). Commit and synchronize. \label{alg:ws_only_gemm1}
\STATE Preserve the previous value $m_i^{\mathrm{old}} = m_i$ and update $m_i = \mathrm{max}(m_i^{\mathrm{old}}, \mathrm{rowmax}(\mathbf{S}_i^{(j)}))$. \label{code-ws:softmax_start}
\STATE Evaluate $\widetilde{\mathbf{P}}_i^{(j)} = \mathrm{exp}(\mathbf{S}_i^{(j)} - m_i)$ and modify $\ell_i = \mathrm{exp}(m_i^{\mathrm{old}} - m_i) \ell_i + \mathrm{rowsum}(\widetilde{\mathbf{P}}_i^{(j)})$. \label{code-ws:softmax_end}
\STATE Wait until $\mathbf{V}_j$ has reached shared memory.
\STATE Update $\mathbf{O}_i$ via: $\mathbf{O}_i = \mathrm{diag}(\mathrm{exp}(m_i^{\mathrm{old}} - m_i))^{-1} \mathbf{O}_i + \widetilde{\mathbf{P}}_i^{(j)} \mathbf{V}_j$ (RS-GEMM). Commit and block for synchronization. \label{alg:ws_only_gemm2}
\STATE Free up the $(j\,\%\,s)$th buffer stage for the producer.
\ENDFOR
\STATE Finalize: compute $\mathbf{O}_i = \mathrm{diag}(\ell_i)^{-1} \mathbf{O}_i$ and $L_i = m_i + \log(\ell_i)$.
\STATE Store the computed $\mathbf{O}_i$ and $L_i$ back to HBM as the $i$th segment of $\mathbf{O}$ and $L$.
\ENDIF
\end{algorithmic}
\end{algorithm}

In our realization of \cref{alg:flash3_wgmma_ws_only} targeting Hopper, we apply \verb|setmaxnreg| to manage register (de)allocations, utilize TMA instructions to fetch $\mathbf{Q}_i$ and $\{\mathbf{K}_j, \mathbf{V}_j\}_{0 \leq j < T_c}$, and employ WGMMA for carrying out the GEMM operations within the mainloop of the consumer, distinguishing the data source of the first operand—shared memory or register file—through the use of the SS or RS prefixes, respectively. To understand the behavior of \cref{alg:flash3_wgmma_ws_only}, one should recognize that TMA loads are issued asynchronously, allowing subsequent loads without waiting for previous ones to finish. Additionally, within the producer’s mainloop, the initial $s$ iterations proceed without any explicit waits since these cycles are dedicated to populating the buffer.

\paragraph{Pingpong scheduling} Due to the asynchronous execution of WGMMA and TMA, as well as the employment of warp-specialization, it becomes feasible to interleave the softmax computation being carried out by one warpgroup with the GEMM operation executed by a different warpgroup. This is particularly relevant considering the significant disparity in throughput between matmul and non-matmul operations on contemporary hardware accelerators. For instance, the H100 SXM5 GPU provides 989 TFLOPS of FP16 matmul compute, yet only achieves 3.9 TFLOPS for executing special functions such as exponentials\footnote{The CUDA programming guide reports that each streaming multiprocessor (SM) can issue 16 special function operations per clock cycle. To arrive at 3.9 TFLOPS, we multiply 16 by the 132 SMs and then by the 1830 MHz clock rate.}, which are critical for softmax calculations. During an FP16 attention forward pass with a head dimension of 128, the number of matmul floating-point operations required exceeds that of exponentials by a factor of 512. Nevertheless, given that exponential operations have a throughput that is 256 times lower than matmuls, their execution time can consume up to half of the total compute cycles, relative to matmul. This imbalance is further amplified in FP8 settings, where matmul throughput increases by 2x, yet the performance of the exponential computations does not improve.

Because the exponential computation is handled by a distinct hardware component—the multi-function unit—it is preferable to align its execution with the matmul operations occurring on the Tensor Cores. To achieve this coordination, we employ synchronization barriers (\texttt{bar.sync} instructions) that enforce the execution order such that the GEMMs (specifically, GEMM1—$\mathbf{P} \mathbf{V}$ for a current iteration, and GEMM0—$\mathbf{Q} \mathbf{K}^\top$ for the subsequent iteration) assigned to warpgroup 1 are completed before those assigned to warpgroup 2. This design allows the softmax computation in warpgroup 1 to overlap with the matmul operations in warpgroup 2. The process then alternates: warpgroup 2 transitions to performing softmax while warpgroup 1 proceeds with GEMMs, in a manner referred to as ``pingpong'' scheduling. An overview of this mechanism can be found in~\cref{fig:pingpong_scheduling}. Although in real implementations this pingpong scheduling rarely achieves perfect overlap as suggested in the illustration, our empirical results indicate a noticeable performance benefit—for example, FP16 forward passes with head size 128 and sequence length 8192 show throughput increasing from 570 TFLOPS to a range of 620–640 TFLOPS.

\paragraph{Attention variants} In the cases of multi-query attention~\citep{shazeer2019fast} and grouped query attention~\citep{ainslie2023gqa}, our method aligns with the strategy outlined in \textsc{FlashAttention-2}, modifying tensor indexing in such a way that redundant storage of $\mathbf{K}$ and $\mathbf{V}$ in HBM is avoided.

\subsection{Intra-warpgroup overlapping GEMMs and softmax}

It is possible to interleave certain instructions from the softmax computation with those from the GEMMs within a single warpgroup. Here, we outline a specific method for achieving this overlap.

In the attention algorithm, operations within the inner loop (main loop) have sequential dependencies that impede parallelization within a single iteration. For example, (local) softmax (lines \ref{code-ws:softmax_start} to \ref{code-ws:softmax_end}) relies on the output $\mathbf{S}_i^{(j)}$ of the first GEMM, while the second GEMM takes its result $\widetilde{\mathbf{P}}_i^{(j)}$ as an operand. Indeed, the wait statements in lines \ref{alg:ws_only_gemm1} and \ref{alg:ws_only_gemm2} of \cref{alg:flash3_wgmma_ws_only} serialize the execution of softmax and GEMMs. However, we can break these dependencies by pipelining across iterations through additional buffers in registers. Pursuing this idea, we propose the following two-stage\footnote{Note that the number of stages of the overlapping scheme is bounded by, but need not equal, the number $s$ of stages in the circular SMEM buffer.} GEMM-softmax pipelining algorithm:

\begin{algorithm}[H]
    \caption{\small\label{alg:flash3_wgmma}\textsc{FlashAttention-3} Consumer Warpgroup Forward Computation}
    \begin{algorithmic}[1]
      \REQUIRE  Input matrices $\mathbf{Q}_i \in \mathbb{R}^{B_r \times d}$ and $\mathbf{K}, \mathbf{V} \in \mathbb{R}^{N \times d}$ reside in HBM; the key block size is $B_c$ so that $T_c = \lceil \frac{N}{B_c} \rceil$.
      \STATE Dynamically assign a fixed quantity of registers in relation to the active number of consumer warps.
      \STATE In on-chip memory, set $\mathbf{O}_i = (0) \in \mathbb{R}^{B_r \times d}$ and initialize $\ell_i$ and $m_i$ as $(0)$ and $(-\infty)$ in $\mathbb{R}^{B_r}$, respectively.
      \STATE Wait for $\mathbf{Q}_i$ and $\mathbf{K}_0$ to become available in shared memory.
      \STATE Use WGMMA to perform $\mathbf{S}_{\mathrm{cur}} = \mathbf{Q}_i \mathbf{K}_0^T$. Ensure operation commit and synchronization.
      \STATE Release the buffer segment allocated to the $0$th stage of $\mathbf{K}$.
      \STATE Derive $m_i$, $\tilde{\mathbf{P}}_{\mathrm{cur}}$, and $\ell_{i}$ from $\mathbf{S}_{\mathrm{cur}}$; apply necessary rescaling to $\mathbf{O}_i$.
      \FOR{$1 \le j < T_c - 1$}
        \STATE Hold until $\mathbf{K}_j$ is transferred to shared memory.
        \label{code:mainloop_start}
        \STATE Calculate $\mathbf{S}_{\mathrm{next}} = \mathbf{Q}_i \mathbf{K}_{j}^T$ via WGMMA. Commit without synchronization. \label{code:first_wgmma}
        \STATE Wait for $\mathbf{V}_{j-1}$ to become ready in shared memory.
        \STATE Update $\mathbf{O}_i$ using $\mathbf{O}_{i} + \tilde{\mathbf{P}}_{\mathrm{cur}} \mathbf{V}_{j-1}$ with WGMMA; commit but defer waiting. \label{code:second_wgmma}
        \STATE Await the completion of WGMMA operation on $\mathbf{Q}_i \mathbf{K}_{j}^T$.
        \STATE Based on $\mathbf{S}_{\mathrm{next}}$, update $m_i$, $\tilde{\mathbf{P}}_{\mathrm{next}}$, and $\ell_{i}$. \label{code:softmax}
        \STATE Ensure completion of the WGMMA $\tilde{\mathbf{P}}_{\mathrm{cur}} \mathbf{V}_{j-1}$, then rescale $\mathbf{O}_i$.
        \STATE Free up the $(j\,\%\,s)$th and $(j-1\,\%\,s)$th buffer stages for $\mathbf{K}$ and $\mathbf{V}$, respectively.
        \STATE Assign $\mathbf{S}_{\mathrm{cur}} \gets \mathbf{S}_{\mathrm{next}}$.
        \label{code:mainloop_end}
      \ENDFOR \STATE Wait for shared memory access to $\mathbf{V}_{T_c - 1}$ to complete.
      \STATE Compute final sum as $\mathbf{O}_{i} = \mathbf{O}_{i} + \tilde{\mathbf{P}}_{\mathrm{last}} \mathbf{V}_{T_c - 1}$, using WGMMA. Commit operation and wait for finish.

\STATE Epilogue: Adjust $\mathbf{O}_{i}$ using $m_{i}$. Derive $L_{i}$ utilizing both $m_{i}$ and $\ell_{i}$. Store $\mathbf{O}_{i}$ and $L_{i}$ in HBM, placing them into the $i$-th segments of $\mathbf{O}$ and $L$, respectively.
\end{algorithmic}
\end{algorithm}

\cref{alg:flash3_wgmma} serves to substitute the consumer path detailed in \cref{alg:flash3_wgmma_ws_only}, thereby forming the full \textsc{FlashAttention-3} algorithm in FP16 precision. Conceptually, WGMMA is utilized interchangeably with asynchronous GEMM. In the primary loop, specified by lines \ref{code:mainloop_start} through \ref{code:mainloop_end}, the second WGMMA computation within iteration $j$ (see line \ref{code:second_wgmma}) is scheduled concurrently with the softmax computations for the subsequent iteration $j+1$ (refer to line \ref{code:softmax}).

Although the pipelined architecture depicted above provides potential performance improvements, several real-world factors must be addressed:
\paragraph{Compiler reordering} Although the pseudocode lays out a sequential, idealized execution of operations, the NVCC compiler routinely reorders instructions to improve efficiency. Such instruction reordering can disturb the intentionally designed alternation between WGMMA and non-WGMMA operations, which may cause either inconsistencies in pipeline execution or reduction in intended performance enhancements. Nevertheless, inspecting the SASS output indicates that the compiler tends to preserve the planned overlap, as explored in Section~\ref{sec:2-stage-sass}.
\paragraph{Register pressure} Avoiding register spilling is important for sustaining peak performance. Nevertheless, the use of a 2-stage pipeline means additional registers are needed for keeping intermediate values and maintaining pipeline context. Specifically, one must hold an additional $\mathbf{S}_{\mathrm{next}}$ in registers, leading to increased usage quantified by $B_r \times B_c \times \text{sizeof}(\text{float})$ per threadblock. This heightened register requirement might restrict the possible block sizes, as enlarging block sizes—a common tuning strategy—also requires more registers. Consequently, practitioners must strike a balance based on performance measurements.
\paragraph{3-stage pipelining} Building on the prior 2-stage design, a 3-stage scheme is proposed, with the goal of overlapping the second WGMMA phase alongside the softmax computation. While this could further exploit Tensor Core resources and elevate utilization, it comes at the cost of demanding even greater register capacity to accommodate the extra pipelining step. As a result, finding the right equilibrium between tile granularity and pipelining depth becomes increasingly complex. More information about this 3-stage pipeline, along with its experimental analysis, is provided in~\cref{sec:3-stage}.

\subsection{Low-precision with FP8}
\label{sec:algofp8}

\textbf{Efficiency: layout transformations.} Executing the forward pass of \textsc{FlashAttention-3} using FP8 precision introduces further difficulties with respect to maintaining layout consistency, issues that are not present when using FP16.

To begin with, it is important to observe that the input tensors $\mathbf{Q}$, $\mathbf{K}$, and $\mathbf{V}$ are commonly stored such that the head dimension is contiguous. However, meeting the k-major requirement of FP8 WGMMA in the second GEMM mandates that $\mathbf{V}$—specifically, the portions of $\mathbf{V}$ loaded into SMEM—should be contiguous along the sequence length dimension. Given that a TMA load operation cannot alter which dimension is contiguous, this imposes two possible approaches: (1) perform a pre-processing transpose of $\mathbf{V}$ in GMEM, or (2) execute an in-kernel transpose of the relevant tiles of $\mathbf{V}$ within SMEM after loading. For the first approach, (1), there are two viable sub-options: (1a) fusing the transpose operation into the epilogue of an earlier process (such as during rotary embedding), or (1b) invoking a dedicated pre-processing transpose kernel\footnote{A well-optimized transpose kernel can run at speeds close to the hardware's bandwidth limit~\citep{colfax_cutlass_transpose_2024}.} to swap the strides between the sequence length and head axes. Despite this, (1a) presents integration challenges when using standard libraries, and (1b) proves too inefficient for memory-limited scenarios like inference.

Instead, for FP8 \textsc{FlashAttention-3}, we choose option (2). Within the kernel, we utilize the LDSM (\verb|ldmatrix|) and STSM (\verb|stmatrix|) instructions, which enable a warp of threads to jointly transfer data between SMEM and RMEM at blocks of 128 bytes.\footnote{As detailed in the PTX specification, LDSM/STSM typically operate on $8 \times 8$ matrices with 16-bit elements \cite[\S 9.7.13.4.15-16]{ptx}. For FP8 support, we achieve this by packing two 8-bit entries together; however, the transpose modes of LDSM/STSM lack the capability to unpack these combined 8-bit values. Consequently, intermediate register shuffling is necessary between the LDSM and STSM instructions to realize the intended tile-level transpose, although we omit further specifics here.} These LDSM/STSM instructions are not only efficient in their use of registers, making it practical to run them in the producer warpgroup, but also can perform layout transpositions simultaneously during memory transfers. In addition, once the initial iteration is complete, it is possible to pipeline the transpose of the next $\mathbf{V}$ tile so that it overlaps with the execution of the two WGMMAs that correspond to the previous $\mathbf{V}$ tile and the current $\mathbf{K}$ tile.

Second, we note that, in contrast to FP16, the FP32 accumulator employed in an FP8 WGMMA possesses a register memory organization that does not align with that of operand A as stored in registers. Illustrations of segments of these distinct layouts are provided in \cref{fig:rmem_accum} and \cref{fig:rmem_operand}, where register allocation per thread follows the sequence indicated. Through the use of byte permutation instructions, we can convert the accumulator output from the first WGMMA to an arrangement compatible with the input requirements of the subsequent WGMMA, as well as with the structure of the $\mathbf{V}$ tile generated via the in-kernel transpose. In detail, and referencing \cref{fig:rmem_accum}, the sequential order is rearranged to
$$\{ \verb|d0 d1 d4 d5 d2 d3 d6 d7| \},$$
and this register ordering is applied repeatedly every 8 bytes. Functionally, this operation permutes the columns of the logical $\mathbf{P}$ tile (e.g., columns $0189$ become the first four). To ensure WGMMA then computes the intended output tile, we can adapt the in-kernel transpose so that it produces a compatible row reordering of the $\mathbf{V}$ tile.\footnote{The flexibility introduced by performing the transpose inside the kernel obviates the need for shuffle operations that transfer register ownership among threads, as described previously in~\citep{colfax_fp8_flashattention_2024}.}

\textbf{Accuracy: block quantization and incoherent processing.} In the FP8 (e4m3) scheme, only 3 bits are devoted to the mantissa while 4 bits represent the exponent. This provides less precision and introduces more numerical errors compared to FP16/BF16. Additionally, large-scale models often exhibit extreme outlier values~\citep{dettmers2208llm,
  sun2024massive}, which far exceed the magnitude of most elements, thereby posing substantial challenges for quantization. Per-tensor scaling~\citep{micikevicius2022fp8} is commonly adopted, where a single scale factor is assigned to each tensor (such as $\mathbf{Q}$, $\mathbf{K}$, and $\mathbf{V}$).
To address the elevated numerical errors in FP8-based attention, we utilize two distinct strategies:

\begin{enumerate}

\item \textbf{Block quantization}: In this method, a single scalar is maintained for each block; consequently, for every tensor in $\mathbf{Q}$, $\mathbf{K}$, and $\mathbf{V}$, we partition them into blocks of dimensions $B_r \times d$ or $B_c \times d$, and quantize the blocks independently. This quantization approach can be combined with operations immediately preceding attention—such as rotary embedding—without incurring extra latency, as such operations are limited by memory bandwidth rather than computation. Since \textsc{FlashAttention-3} natively processes data in blocks, each block within $\mathbf{S}$ can be rescaled to compensate for block-wise quantization, all without introducing additional computational overhead.
\item \textbf{Incoherent processing}: To mitigate the effects of outliers, we apply a random orthogonal matrix $\mathbf{M}$ to both $\mathbf{Q}$ and $\mathbf{K}$ before executing FP8 quantization. Owing to the orthogonality property of $\mathbf{M}$, where $\mathbf{M} \mathbf{M}^\top = I$, it follows that $(\mathbf{Q} \mathbf{M}) (\mathbf{K} \mathbf{M})^\top = \mathbf{Q} \mathbf{K}^\top$, which ensures the attention computation remains unchanged when both $\mathbf{Q}$ and $\mathbf{K}$ are transformed by $\mathbf{M}$. This process effectively redistributes the influence of outliers, as each component of $\mathbf{Q} \mathbf{M}$ or $\mathbf{K} \mathbf{M}$ comprises a random linear combination of elements from the original tensor, leading to decreased quantization errors. Consistent with the methodology of \citet{chee2024quip} and~\citet{tseng2024quip}, $\mathbf{M}$ is constructed as a product of random diagonal sign matrices (entries $\pm 1$) and a Hadamard matrix, thus supporting multiplication in $O(d \log d)$ time complexity instead of $O(d^2)$, and allowing for integration with rotary embedding operations at no additional computational expense.
\end{enumerate}

Our evaluation in \cref{sec:numerical_error} confirms that these strategies can lower numerical error by as much as a factor of 2.6.

\section{Empirical Validation}
\label{sec:exp}

Our implementation of \textsc{FlashAttention-3} leverages CUTLASS~\citep{Thakkar_CUTLASS_2023} primitives, specifically incorporating WGMMA and TMA abstractions, to assess both its performance and precision.

\begin{itemize}

\item \textbf{Benchmarking attention.} We assess the performance of \textsc{FlashAttention-3} over various sequence lengths and benchmark it against the standard PyTorch implementation, \textsc{FlashAttention-2}, the Triton version of \textsc{FlashAttention-2} optimized with H100-specific instructions, as well as a vendor-supplied, cuDNN-tuned implementation of \textsc{FlashAttention-2} designed for H100 GPUs. Our experiments demonstrate that \textsc{FlashAttention-3} achieves up to a 2.0$\times$ speedup over \textsc{FlashAttention-2} and is 1.5$\times$ faster compared to the Triton-based \textsc{FlashAttention-2}. In terms of computational throughput, \textsc{FlashAttention-3} attains up to 740 TFLOPs/s, which corresponds to 75\% of the peak TFLOPs/s on H100 GPUs.
\item \textbf{Ablation study.} We show through ablation experiments that key algorithmic advances, specifically warp-specialization and the GEMM-softmax pipeline, are substantial contributors to the increased efficiency of \textsc{FlashAttention-3}.
\item \textbf{Accuracy of FP8 attention.} We empirically demonstrate that incorporating block quantization and incoherent processing techniques leads to a 2.6$\times$ reduction in numerical error for FP8 \textsc{FlashAttention-3}.
\end{itemize}

\subsection{Benchmarking Attention}
\label{subsec:benchmark_attn}

We evaluate the execution time of various attention mechanisms on an H100 80GB SXM5 GPU across multiple configurations, considering both the presence and absence of a causal mask, and head dimensions set to 64 or 128, utilizing FP16 inputs. The performance outcomes are presented in~\cref{fig:benchmark_attn_fwd} and~\cref{fig:benchmark_attn_bwd}. Our measurements indicate that \textsc{FlashAttention-3} achieves forward pass speeds approximately 1.5-2.0$\times$ faster than those of \textsc{FlashAttention-2}, and backward pass speeds 1.5-1.75$\times$ higher. Relative to the conventional attention implementation, \textsc{FlashAttention-3} attains improvements of up to 3-16$\times$ in speed. Notably, for sequence lengths of 1,000 tokens and greater, \textsc{FlashAttention-3} even exceeds the throughput of a proprietary vendor library (cuDNN -- closed source) that is specifically optimized for H100 GPUs.

\paragraph{Benchmark settings:} We vary the sequence length as 512, 1k, ..., 16k, and set batch size so that the total number of tokens is 16k. We set the hidden dimension to 2048, and head dimension to be either 64, 128, or 256 (i.e., 32 heads, 16 heads, or 8 heads). To calculate the FLOPs of the forward pass, we use:

\begin{equation*}
  4 \cdot \text{seqlen}^2 \cdot \text{head dimension} \cdot \text{number of heads}.
\end{equation*}

Applying causal masking leads us to halve this figure, reflecting that roughly fifty percent of the elements are actually computed. For the backward pass, we estimate FLOPs by scaling the forward pass value by a factor of 2.5, as the backward computation involves five matrix multiplications in contrast to two during the forward operation, considering recomputation.

Additionally, we assess the forward pass runtime using FP8 under comparable conditions. The outcomes for a headdim of 256 are presented in ~\cref{fig:benchmark_attn_fp8}, with comprehensive results detailed in \cref{sec:benchmark_attn_fp8_full}.

\subsection{Ablation Study: 2-Stage Pipelining Experiments}
\label{sec:2-stage-pipelining-experiments}

We evaluate the impact of removing both the 2-stage WGMMA-softmax pipelining and warp-specialization features in the non-causal FP16 \textsc{FlashAttention-3}, using constant parameters $\{\text{batch}, \text{seqlen}, \text{nheads}, \text{hdim}\} = \{ 4, 8448, 16, 128\}$. The data presented in~\cref{table:ablation_pipelining} demonstrates that the integration of our algorithmic enhancements—specifically, asynchrony via warp-specialization and the overlap of GEMM and softmax operations—produces a substantial increase in speed, improving performance from 570 up to 661 TFLOPs.

\subsection{Numerical Error Validation}
\label{sec:numerical_error}

Given the growing attention to numerical error in \textsc{FlashAttention}~\citep{golden2024flash}, we evaluate the numerical behavior of \textsc{FlashAttention-2}, \textsc{FlashAttention-3}, and a conventional attention implementation by comparing each against a high-precision FP64 reference. To emulate the presence of outlier features and activations characteristic of LLMs~\citep{dettmers2208llm, sun2024massive}, we sample the elements of $\mathbf{Q}, \mathbf{K}, \mathbf{V}$ from the following distribution:

\begin{equation*}
  \mathcal{N}(0, 1) + \mathcal{N}(0, 100) \cdot \mathrm{Bernoulli}(0.001).
\end{equation*}

In this setup, all entries follow a normal distribution with mean zero and standard deviation 1, except that an additional independent component, itself normally distributed with a standard deviation of 10, is added to 0.1\% of the entries. The root mean squared error (RMSE) is then computed as summarized in~\cref{table:numerical_error}. When using FP16 precision, both \textsc{FlashAttention-2} and \textsc{FlashAttention-3} show a 1.7$\times$ reduction in RMSE relative to the default implementation, owing to the fact that intermediate values (softmax outputs) are preserved in FP32. The baseline FP8 implementation relies on per-tensor scaling, with matrix multiplication accumulators in FP32 and intermediate softmax steps in FP16. Leveraging block quantization and incoherent computation, \textsc{FlashAttention-3} in FP8 attains 2.6$\times$ better accuracy than the baseline.

\section{Dicussion, Limitations, Conclusion}
\label{sec:discussion}

Through the development of \textsc{FlashAttention-3}, we have shown that advances in programming methods and the adoption of hardware capabilities such as asynchrony and low-precision computation can significantly improve both the performance and fidelity of attention mechanisms. Our approach achieves a 1.5-2.0$\times$ increase in attention throughput relative to \textsc{FlashAttention-2}, and it reduces FP8 numerical errors by a factor of 2.6 compared to conventional per-tensor quantization. Nonetheless, several aspects of our implementation remain to be addressed, including further optimizing for LLM inference workloads, incorporating a persistent kernel into the FP8 variant,\footnote{In our performance evaluations, the FP16 version of \textsc{FlashAttention-3} leverages a persistent kernel as well as a load balancing strategy, while this is currently not available in the FP8 variant. This partially accounts for FP8 \textsc{FlashAttention-3} underperforming with short sequence lengths and causal masking in comparison to the FP8 cuDNN kernels.} and systematically evaluating the implications of low-precision attention when training at scale. While our experiments have centered on Hopper GPUs, we anticipate the methodologies introduced here are applicable to a broader range of hardware accelerators. Ultimately, we envision that improvements to core operators such as attention can facilitate new directions in long-context modeling and applications.

\end{document}